\documentclass[11pt]{article}
\usepackage[a4paper,margin=1in]{geometry}
\usepackage{amsmath,amssymb,amsthm,booktabs,microtype}
\newtheorem{theorem}{Theorem}[section]
\newtheorem{proposition}[theorem]{Proposition}
\theoremstyle{remark}\newtheorem{remark}[theorem]{Remark}
\title{The Sharp Logarithmic Block Clock for Hilbert--Schmidt Spectral Tails}
\author{Bin Wang}\date{July 2026}
\begin{document}\maketitle
\begin{abstract}
We determine the exact logarithmic clock forced by a Hilbert--Schmidt mass
law.  Let a normal diagonal spectral tail have eigenvalue moduli at most $q$
and squared mass $M_\sigma\le C\sigma^{-\alpha}$.  On a disk of radius $R$
with $qR<1$, a clock $m_\sigma=\lceil a\log(1/\sigma)\rceil$ suppresses the
high-order regularized-determinant tail when
$a\log(1/(qR))\ge\alpha$.  Equality gives only a logarithmic gain.  Below
that boundary a diagonal saturation family makes the tail diverge, so the
threshold is sharp for this information class.  The mass-and-cap root-rate
upper bound is strictly below one under strict inequality, but a particular
nonsaturating family may satisfy the RH-279 test at a smaller slope.
\end{abstract}

\section{Mass-and-cap class}
Let $C_\sigma=\operatorname{diag}(\mu_j(\sigma))$ with
$|\mu_j(\sigma)|\le q$ and
\[
 M_\sigma=\sum_j|\mu_j(\sigma)|^2\le C\sigma^{-\alpha},
 \qquad \alpha>0.
\]
For $qR<1$ and $m\ge2$, the same estimate as in RH-282 gives
\begin{equation}\label{eq:tail}
 T_\sigma(m,R):=
 \sum_{n\ge m}\frac{|\operatorname{Tr}C_\sigma^n|R^n}{n}
 \le\frac{Cq^{-2}\sigma^{-\alpha}(qR)^m}
 {m(1-qR)}.
\end{equation}

\section{Sharp clock theorem}
\begin{theorem}\label{thm:clock}
Put $m_\sigma=\lceil a\log(1/\sigma)\rceil$.  If
\[
 a\log\frac1{qR}>\alpha,
\]
then $T_\sigma(m_\sigma,R)$ decays by a positive power of $\sigma$.  If
equality holds, the right side of \eqref{eq:tail} is
$O(1/\log(1/\sigma))$ and still tends to zero.

If $a\log(1/(qR))<\alpha$, there exists a diagonal family satisfying the
same mass and cap hypotheses for which $T_\sigma(m_\sigma,R)\to\infty$.
\end{theorem}
\begin{proof}
The upper assertions follow after writing
$(qR)^{m_\sigma}=O(\sigma^{a\log(1/(qR))})$ in \eqref{eq:tail}.
For sharpness take
$N_\sigma=\lfloor c\sigma^{-\alpha}\rfloor$ copies of $q$, with $c>0$
small enough that $N_\sigma q^2\le C\sigma^{-\alpha}$.  Then
\[
 T_\sigma(m,R)=N_\sigma\sum_{n\ge m}\frac{(qR)^n}{n}.
\]
The first term gives
$T_\sigma\ge N_\sigma(qR)^m/m$, which diverges when the displayed strict
reverse inequality holds.
\end{proof}

\begin{proposition}[uniform root-rate boundary]
For the saturation scale $M_\sigma\asymp\sigma^{-\alpha}$,
\[
 \limsup_{\sigma\downarrow0}
 \|C_\sigma^{m_\sigma}\|_1^{1/m_\sigma}R
 \le qR e^{\alpha/a}.
\]
The right-hand side is strictly smaller than one exactly when
$a>\alpha/\log(1/(qR))$.  Thus the mass-and-cap information uniformly
guarantees the RH-279 root-rate test in the strict supercritical region.
For the repeated-$q$ saturation family used above, the actual root-rate
limit equals $qRe^{\alpha/a}$, so this boundary is exact for that family.
\end{proposition}
\begin{proof}
Use $\|C^m\|_1\le M_\sigma q^{m-2}$ and take $m$th roots.  For
$N_\sigma\asymp\sigma^{-\alpha}$ repeated entries equal to $q$,
$\|C_\sigma^{m_\sigma}\|_1=N_\sigma q^{m_\sigma}$, which gives equality
in the limiting root rate.
\end{proof}

\section{Hardy-disk instance and boundary}
For $\alpha=1$, $q=1/2$, and $R=7/5$,
\[
 a_{\rm crit}=\frac1{\log(10/7)}=2.803673252\ldots .
\]
The RH-282 choice $a=4$ has strict exponent
$4\log(10/7)-1=0.426699775\ldots$.

The lower model proves sharpness only from mass and modulus information.  It
does not claim that the folded Gaussian spectrum has repeated roots at $q$.
The spectral-to-monodromy bridge and Gates A--E remain open.
\end{document}
